\scanpage{17}
\mspage{- IV-\handwritten{962} -}
que de $k$ ; il est alors immédiat que $K\otimes_k\Omega$ est composé direct d'un nombre
fini de corps , tous isomorphes à la clôture séparable $K\Omega$ de $K$ . D'autre
part , l'anneau $ {}^{\mathrm{hs}}A$ étant hensélien , $C\otimes_A({}^{\mathrm{hs}}A)$ est composé direct des annea\handwritten{ux}
locaux $(C\otimes_A({}^{\mathrm{hs}}A))_{\underline{q}_i}$ , où $\underline{q}_i$ $(1\leq i\leq s)$ parcourt les idéaux maximaux de \deleted{\illegible{1}}
$C\otimes_A({}^{\mathrm{hs}}A)$ , et ces anneaux locaux sont henséliens (18.5.10) et ont $K\Omega$ pour
corps résiduel , donc sont strictement locaux (18.8.1) . Posons pour abréger
$C'=C\otimes_A({}^{\mathrm{hs}}A)$ et $C'_i=C'_{\underline{q}_i}$ , qui est donc un anneau quotient de $C'$ ; soient
$f:C\to C'$ , $g:{}^{\mathrm{hs}}A\to C'$ , $\handwritten{\varphi}_i:C'\to C'_i$ les applications canoniques .

\diagram{\handwritten{\begin{tikzcd}[ampersand replacement=\&,column sep=small,row sep=small]
\& \& \& \handwritten{{}^{\mathrm{hs}}C}\arrow[dl,bend right=10,"\handwritten{w'}"']\\
\handwritten{C}\arrow[r,"\handwritten{f}"']\arrow[urrr,"\handwritten{v}"] \& \handwritten{C'}\arrow[r] \& \handwritten{C'_i}\arrow[ur,bend right=10,"\handwritten{w}"'] \& \\
\handwritten{A}\arrow[u]\arrow[r,"\unsure{\prime|/}"] \& \handwritten{{}^{\mathrm{hs}}A}\arrow[u,-]\arrow[uurr,"\handwritten{u}"'] \& \&
\end{tikzcd}}}

Soit $v:$ \deleted{\illegible{1}} \inserted{\handwritten{$C\to{}^{\mathrm{hs}}C$}} l'homomorphisme canonique ; il existe \deleted{\illegible{1}} un homomorphisme
local et un seul $u:{}^{\mathrm{hs}}A\to{}^{\mathrm{hs}}C$ qui , par passage aux quotients,donne l'injection
canonique $\handwritten{\Omega}\to K\Omega$ \deleted{\illegible{1}} et tel que le composé $A\longrightarrow C\xrightarrow{v}{}^{\mathrm{hs}}C$
soit égal au composé $A\longrightarrow{}^{\mathrm{hs}}A\xrightarrow{u}{}^{\mathrm{hs}}C$ (18.8.5) ; donc il existe un unique
homomorphisme local $w_{\unsure{o|0}}:C'\longrightarrow{}^{\mathrm{hs}}C$ tel que $w_o\handwritten{\circ g=u}$, $w_o\handwritten{\circ f=v}$, et puisque $ {}^{\mathrm{hs}}C$
est un anneau local , cet homomorphisme se factorise en $C'\xrightarrow{\handwritten{\varphi_i}}C'_i\xrightarrow{w}{}^{\mathrm{hs}}C$
pour un indice $i$ bien déterminé, $w$ étant un homomorphisme local . D'autre
part , l'anneau $C'_i$ étant strictement local , il existe un homomorphisme local
$w':{}^{\mathrm{hs}}C\to C'_i$ qui , par passage aux quotients , donne l'automorphisme
identique de $K\Omega$ et qui est tel que $w'\circ v=\handwritten{\varphi}_i\circ f$ \inserted{(18.8.5)} .\deleted{\illegible{1}} On déduit d'abord de là
que $w\circ w'$ est un endomorphisme de $ {}^{\mathrm{hs}}C$ qui , par passage aux quotients \deleted{\unsure{à|a}},donne
l'automorphisme identique de $K\Omega$, donc est l'identité (18.8.5) . D'autre
part , , on a $w'\circ w\circ\handwritten{\varphi}_i\circ f=w'\circ v=\handwritten{\varphi}_i\circ f$ , et $w'\circ w\circ\handwritten{\varphi}_i\circ g=w'\circ u=\handwritten{\varphi}_i\circ g$ ; d'où $w'\circ w\circ\handwritten{\varphi}_i=
\handwritten{\varphi}_i$ , autrement dit $w'\circ w$ est l'automorphisme identique de $C'_i$ . CQFD .

\uline{Remarque} (18.8.{\deleted{7}\handwritten{10}}.1) .- La démonstration de (18.8.7)
% slanted
\handwritten{\unsure{!}}\inserted{\handwritten{! pour une $A$-algèbre $B$ \uline{entière} sur $A$}} n'ut\uline{i}lise le
fait que $B$ est une $A$-algèbre finie,que pour établir que $K\otimes_k\Omega$ est
composé direct d'un nombre fini de corps . Or , cette dernière propr\handwritten{i}été
est encore vérifiée lorsque l'on suppose que $B$ est une $A$-algèbre
\uline{entière semi-locale} dont les idéaux maximaux $\underline{n}_j$ ont des corps résiduels
de degrés séparables $\handwritten{[\kappa}(\underline{n}_j):k\handwritten{]}_s$ \uline{finis} .

\uline{Proposition} (18.8.\deleted{8}\handwritten{11}) .- \uline{Soit} $A$ \uline{un anneau} \deleted{\illegible{1}} \uline{semi-lo\unsure{d|c}al} .

(i) \uline{Pour que} $ {}^{\mathrm{hs}}A$ \uline{soit réduit} (resp. \uline{normal) , il faut et il suffit que} $A$ \uline{le soit} .

(ii) \uline{Supposons} $A$ \uline{noethérien . Alors , pour tout idéal premier} $\underline{p}$ \uline{de} $A$ ,\uline{l'anneau} $({}^{\mathrm{hs}}A)_{\underline{p}}/\underline{p}.({}^{\mathrm{hs}}A)_{\underline{p}}$ \uline{est composé direct d'un nombre fini d'extensions \unsure{algé|algè}-}
% archive sheet number: 16
